% Chapter 11: calibration.
\chapter{Calibration}\label{ch:calibration}
\index{calibration}

Calibration teaches the machine two things: how far the roll \emph{coasts}
after the motor stops, and (optionally) how long an inch really is on your
measuring wheel. It takes about 10 to 15 minutes.

\textbf{Calibrate when:} the machine is new; you have updated the firmware
from a version older than fw=6; the wheel has been replaced, cleaned of
build-up, or re-mounted; the motor, belt or roll core has changed; or rolls
are often much longer than the label.

\section{Before you start}
\begin{itemize}
  \item Connect the PC and start \textbf{winder} (Chapter~\ref{ch:pc}). The
        \ui{Home} tab should show the machine.
  \item Load paper exactly as for a real wind. The machine will run several
        short test winds and use some paper.
  \item Have a tape measure or ruler and a pencil ready for the length check.
  \item Want to see the steps first? Click \ui{Practice (nothing moves)}.
\end{itemize}

\section{The four steps}
On \ui{Home} click \ui{Calibrate the machine}, or open the \ui{Calibrate} tab.

\begin{enumerate}[label=\textbf{Step \arabic*},leftmargin=4.5em,itemsep=0.7em]
  \item \textbf{Get ready.} Read the checklist on screen and confirm.
  \item \textbf{Test winds.} Click \ui{Start automatic test}. The machine
        runs at least four short test winds, measures how far each one
        coasts and works out the average. It then runs a check wind; the
        check passes if it stops within 20 counts (about 0.1 inch) of where
        it was told to. Stay clear of the machine; it starts and stops by
        itself. Click \ui{Stop motor} at the top of the page if anything goes
        wrong.
  \item \textbf{Check the length (optional, recommended).} This corrects the
        wheel size.
        \begin{enumerate}[label=\alph*.]
          \item Draw a pencil line across the paper at the edge of the table.
          \item Click \ui{Run}. The machine winds about 36 inches.
          \item Measure from the table edge to the pencil line with a ruler, or count
                sheets laid along it (US letter is 11 inches long, A4 is
                11.69 inches).
          \item Type in what you measured.
        \end{enumerate}
        Very small differences (under 0.2\%) are ignored. A difference over
        10\% is refused as a measuring mistake; measure again. If the wheel
        size changes, the program runs the automatic test again with the new
        size.
  \item \textbf{Save.} Look over the new numbers and click
        \ui{Save to the machine}. Only the numbers you reviewed are stored.
        Then click \ui{Back to Home}.
\end{enumerate}

If the program shows \ui{Something went wrong}, read the message, then click
\ui{Try again} or \ui{Leave}. Nothing is saved until Step 4.

\section{What calibration changes}
\begin{itemize}
  \item \ui{Coast after the motor stops} (\setting{coast}), from the test winds.
  \item \ui{Measuring wheel size} (\setting{wheel\_um}), from the length check
        if you did it.
  \item It never changes \setting{margin} or \setting{bump}, or the
        target lengths.
\end{itemize}

\section{After calibrating}
Wind one 30 yard roll and look at the \shows{DONE} screen. A few inches over
is normal. As an extra check, measure a roll by hand now and then and keep a
note in the maintenance log (Chapter~\ref{ch:maintenance}).

\section{From the command line}
A technician can also run \texttt{winder calibrate} in a terminal, with
\texttt{--lengths}, \texttt{--measured} and \texttt{--save} (see
\texttt{winder --help}).
